\section{Introduction}\label{sec:introduction}

% Replace the text of every chapter with your own. The paragraphs below only
% show the commands of the template; README.md explains each of them.

\subsection{Background}\label{sec:background}
Describe the context of the problem here. A key concept is introduced with \verb|\term|, for example \term{socio-technical system}. Sources are cited with a non-breaking space before the citation, as in the ACM Code of Ethics~\cite{acm} or a book on technology governance~\cite{collingridge}.

Abbreviations are declared once in \texttt{config/abbreviations.tex}. The first use of \ac{ai} is written out in full; later uses of \ac{ai} print the short form only.

\subsection{Problem Statement}\label{sec:problem}
State the problem the report answers. A highlighted box suits the central question:

\begin{keybox}{Core problem}
  \textit{Write the one-sentence question that the report answers.}
\end{keybox}

\subsection{Objectives and Scope}\label{sec:objectives}
List what the report sets out to do:

\begin{itemize}
  \item \textbf{Objective 1:} describe the first goal.
  \item \textbf{Objective 2:} describe the second goal.
\end{itemize}

\subsection{Report Structure}\label{sec:structure}
Cross-references are written with \verb|\cref| so that their numbers follow the document. \Cref{sec:method} describes the method, \cref{sec:results} presents the results and \cref{sec:conclusion} concludes. Supporting material is in \cref{app:example}.
